\documentclass[11pt]{article}
\usepackage[T1]{fontenc}
\usepackage[margin=1in]{geometry}
\usepackage{amsmath,amssymb}

\setlength{\parindent}{0pt}
\setlength{\parskip}{0.65\baselineskip}

\newcommand{\Type}[1]{\texttt{#1}}
\newcommand{\Field}[1]{\ensuremath{\mathrm{#1}}}

\title{Proposed Mathematical Additions\\to the PuppyCAD Design Graph Specification}
\author{Draft Additions}
\date{July 16, 2026}

\begin{document}
\maketitle
\tableofcontents
\newpage

\section{Purpose and Scope}

The PuppyCAD specification describes a design as a graph containing authored
intent, derived geometry, constraints, and diagnostics. This document proposes
a more precise mathematical foundation for that graph without prescribing an
implementation, programming language, storage format, geometry kernel, or
solver algorithm.

The additions are intended to answer questions such as:

\begin{itemize}
  \item What exactly is the graph at a particular moment?
  \item When is a node well-typed?
  \item What is the difference between a dependency and a geometric relation?
  \item What does a reference resolve to?
  \item What does applying an event mean?
  \item When is evaluator output stale?
  \item What does it mean for a constraint to be satisfied?
  \item Can a graph be valid when a feature has failed?
\end{itemize}

The proposed definitions describe observable semantics. Implementations remain
free to choose their own data structures, serialization, hashing algorithm,
solver, kernel, cache, and user interface as long as those choices preserve the
semantics.

\subsection{Non-goals}

This document does not define:

\begin{itemize}
  \item a textual or binary file format,
  \item a concrete representation of identifiers, references, or argument values,
  \item a particular event storage strategy,
  \item a numerical constraint-solving algorithm,
  \item a solid-modeling or topology kernel,
  \item a query language or selector syntax,
  \item an API, command-line interface, or user interface.
\end{itemize}

These may be specified separately. Their observable behavior should conform to
the abstract semantics defined here.

\section{Common Notation}

Let:

\begin{itemize}
  \item $I$ be the universe of node identifiers,
  \item $T$ be the universe of node types,
  \item $V$ be the universe of argument values,
  \item $A$ be the universe of authorities, such as users, agents, solvers,
    kernels, and importers,
  \item $\mathcal{R}$ be the universe of references,
  \item $\mathcal{P}(X)$ be the power set of $X$, and
  \item $X^*$ be the set of finite ordered sequences over $X$.
\end{itemize}

A node has the existing PuppyCAD form:

\[
n=(\operatorname{id}(n),\operatorname{type}(n),
   \operatorname{src}(n),\operatorname{arg}(n)),
\]

where:

\[
\operatorname{id}(n)\in I,\qquad
\operatorname{type}(n)\in T,\qquad
\operatorname{src}(n)\in\mathcal{R}^*,\qquad
\operatorname{arg}(n)\in V.
\]

The precise representation of a reference or argument value is deliberately
left unspecified.

\section{Graph Snapshots and Event History}

\subsection{Motivation}

The current graph definition places the live node set, root set, and event log
in one tuple. This makes it unclear whether the node set is independently
authoritative or the result of folding the log. It also makes implementations
that retain only a snapshot, only a log, or both difficult to describe.

\subsection{Graph snapshot}

A \textbf{graph snapshot} is:

\[
S=(N,R),
\]

where $N$ is a finite set of live nodes and $R\subseteq I$ is a finite set of
root identifiers. If a domain needs a presentation order among roots, that
order is a separate relation rather than part of root membership.

The following identity invariants hold:

\[
\forall n_1,n_2\in N:\
\operatorname{id}(n_1)=\operatorname{id}(n_2)\implies n_1=n_2,
\]

and:

\[
R\subseteq\operatorname{ids}(N),
\qquad
\operatorname{ids}(N)=\{\operatorname{id}(n)\mid n\in N\}.
\]

\subsection{History and folding}

A \textbf{history} is a finite ordered sequence of events:

\[
H=(e_1,e_2,\ldots,e_k).
\]

Event application is a partial function:

\[
\operatorname{apply}:\mathcal{S}\times E\rightharpoonup\mathcal{S},
\]

where $\mathcal{S}$ is the universe of snapshots. It is partial because an
event may not be applicable to a particular snapshot. Adding a node whose
identifier is already live is one example.

Folding a history is defined recursively:

\[
\operatorname{fold}(S_0,())=S_0
\]

and:

\[
\operatorname{fold}(S_0,(e_1,\ldots,e_k))
=
\operatorname{apply}
\left(
  \operatorname{fold}(S_0,(e_1,\ldots,e_{k-1})),e_k
\right).
\]

The current snapshot is therefore:

\[
S_k=\operatorname{fold}(S_0,H).
\]

An implementation may store $S_k$, $H$, or both. If it stores both, they must
satisfy the consistency invariant above.

\subsection{Example}

Suppose the initial snapshot contains only an XY reference plane:

\[
S_0=(\{p_{xy}\},\{\operatorname{id}(p_{xy})\}).
\]

The history contains three events:

\begin{enumerate}
  \item add sketch $s_1$ on $p_{xy}$,
  \item add rectangle entities to $s_1$, and
  \item add extrusion $x_1$ from a profile in $s_1$.
\end{enumerate}

Then:

\[
S_3=\operatorname{fold}(S_0,(e_1,e_2,e_3)).
\]

Snapshot $S_3$ describes what currently exists. The history explains how that
state arose. These are related but distinct concepts.

\section{The Typed Attributed Graph}

\subsection{Motivation}

Listing node types describes an intended vocabulary, but it does not formally
state which source references and arguments are legal for each type. Without a
typing relation, the specification cannot distinguish a malformed extrusion
from a valid one at the abstract level.

\subsection{Type environment}

Let a \textbf{type environment} be a mapping:

\[
\Gamma:T\rightarrow
(\operatorname{SrcSig},\operatorname{ArgSig},\operatorname{Class}).
\]

For every type $t$, $\Gamma(t)$ describes:

\begin{itemize}
  \item the expected roles and target types of source references,
  \item the mathematical shape and dimensions of its arguments, and
  \item its semantic class, such as authored, derived, or diagnostic.
\end{itemize}

The judgment:

\[
\Gamma;S\vdash n
\]

means that node $n$ is well-typed in snapshot $S$. The judgment:

\[
\Gamma\vdash S
\]

means that the complete snapshot is well-typed.

At minimum, well-typedness requires:

\begin{enumerate}
  \item every live node has a unique identifier,
  \item every root identifies a live node,
  \item every required reference resolves with an accepted cardinality,
  \item resolved targets have types permitted by the source signature,
  \item arguments have the required value kind and physical dimension, and
  \item node-local invariants hold.
\end{enumerate}

\subsection{Example: extrusion}

An abstract type rule for an extrusion may require exactly one profile source,
a distance having dimension $L$, a direction vector or inherited profile
normal, and an operation that creates, joins, or cuts material:

\[
\frac{
  \operatorname{resolve}_S(r_p)=\{p\}
  \quad
  \operatorname{type}(p)=\Type{Profile}
  \quad
  \dim(d)=L
  \quad
  op\in\{\mathrm{create},\mathrm{join},\mathrm{cut}\}
}{
  \Gamma;S\vdash
  (i,\Type{Extrude},(r_p),(d,op))
}.
\]

This rule does not say how an extrusion is represented in memory or how a solid
is calculated. It defines only when the authored fact is meaningful.

\subsection{Open vocabulary}

An open vocabulary is compatible with formal typing. An extension introduces a
new type together with its source signature, argument signature, and semantic
class. Qualified type names are useful at the abstract level because the
existing vocabulary uses names such as \Type{Shell} for more than one concept.
For example:

\[
\Type{feature:Shell}\neq\Type{topology:Shell}.
\]

The exact spelling is not important. The important invariant is that a type
denotes one semantic contract within a type environment.

\section{Edge Roles and Graph Structure}

\subsection{Motivation}

The source field currently carries several meanings. An extrusion consumes a
profile, a perpendicular constraint relates two lines, a diagnostic observes a
failed feature, and a generated face records provenance from an extrusion.
Treating all of these as one undifferentiated dependency makes evaluation and
cycle rules unclear.

\subsection{Labelled edges}

Each source position has a role from a set $K$, for example:

\[
K=\{\mathrm{input},\mathrm{relation},\mathrm{provenance},
     \mathrm{observation},\mathrm{containment}\}.
\]

Resolving source references induces a labelled edge relation:

\[
D_S\subseteq I\times K\times I.
\]

An edge $(u,k,v)\in D_S$ means that node $v$ refers to node $u$ in role $k$.
A useful projection is:

\[
D_S^{\mathrm{input}}
=
\{(u,v)\mid(u,\mathrm{input},v)\in D_S\}.
\]

Corresponding projections exist for provenance, relations, observations, and
containment. The specification may impose different invariants on them:

\begin{itemize}
  \item evaluation inputs may be required to form an acyclic graph,
  \item symmetric constraint relations may contain cycles,
  \item containment may be required to form a forest, and
  \item diagnostics must not become material-producing inputs unless a type
    contract explicitly permits it.
\end{itemize}

\subsection{Example}

Consider two lines constrained to be perpendicular and then used to define a
profile. The constraint edges are relational; the path from profile to
extrusion is an evaluation dependency.

\begin{verbatim}
line-a -- relation --> perpendicular-constraint <-- relation -- line-b
   |                                                               |
   +--------------------- input -----------------------------------+
                                  |
                               profile
                                  |
                                input
                                  |
                               extrude
\end{verbatim}

The constraint subgraph may contain loops. The material evaluation path is
directional. Invalidation algorithms need the latter, while a constraint solver
needs the former.

\section{Reference Semantics}

\subsection{Motivation}

A reference status alone does not state what was selected, whether multiple
results are intentional, or what cardinality its consumer requires. Staleness
is also conceptually different from missing or ambiguous resolution.

\subsection{Set-valued resolution}

A reference resolves to a set of targets:

\[
\operatorname{resolve}_S:
\mathcal{R}\rightarrow\mathcal{P}(\operatorname{Target}(S)).
\]

A target may be a whole node or a semantic subfeature of a node. A reference
also has an expected cardinality:

\[
\operatorname{card}(r)\in
\{\mathrm{one},\mathrm{zeroOrOne},\mathrm{oneOrMore},\mathrm{many}\}.
\]

For a reference that expects exactly one target:

\[
\operatorname{resolutionStatus}_S(r)=
\begin{cases}
\mathrm{missing} & |\operatorname{resolve}_S(r)|=0,\\
\mathrm{valid} & |\operatorname{resolve}_S(r)|=1,\\
\mathrm{ambiguous} & |\operatorname{resolve}_S(r)|>1.
\end{cases}
\]

For a set-valued query, multiple results may be valid. Ambiguity is therefore
not equivalent to resolving more than one target. It means resolving to a
cardinality not accepted by the consuming source position.

Staleness is not a reference-resolution status. It is a relationship between a
derived fact and the current versions of its sources.

\subsection{Example: semantic face selector}

Suppose an extrusion produces two planar cap faces and several side faces. A
selector requests ``the cap face whose outward normal is $+Z$.'' If exactly one
face satisfies that semantic predicate, the selector is valid.

If a later operation splits the cap into two disconnected coplanar faces, the
same selector may resolve to two faces. A consumer expecting one face now sees
\Type{ambiguous}; a consumer accepting a face set may continue to be valid.

\subsection{Stability expectation}

The specification need not prescribe a topology-naming algorithm, but it should
define the desired invariance. If snapshots $S$ and $S'$ preserve the semantic
evidence named by selector $r$, then resolution should preserve the
corresponding target under a topology correspondence $\phi$:

\[
\operatorname{resolve}_{S'}(r)
=
\phi(\operatorname{resolve}_S(r)).
\]

When no unique correspondence exists, the result should be explicitly missing
or ambiguous rather than silently rebound by creation order.

\section{Events as Atomic State Transitions}

\subsection{Motivation}

The rewrite vocabulary states the intended mutations but not their
preconditions, atomicity, or failure behavior.

\subsection{Event transition}

Write:

\[
S\xrightarrow{e}S'
\]

when applying event $e$ to snapshot $S$ produces snapshot $S'$. An event
contains an ordered finite sequence of rewrites:

\[
e=(i,a,(w_1,\ldots,w_m),P),
\]

where $i$ is the event identity, $a\in A$ is its authority, and $P$ is its
provenance or precondition record.

An event is \textbf{atomic}: either every rewrite is applicable in order and the
complete transition occurs, or the snapshot remains unchanged.

\subsection{Add node}

For $S=(N,R)$:

\[
\frac{
  \operatorname{id}(n)\notin\operatorname{ids}(N)
  \quad
  \Gamma;(N\cup\{n\},R)\vdash n
}{
  (N,R)\xrightarrow{\operatorname{AddNode}(n)}(N\cup\{n\},R)
}.
\]

\subsection{Replace node}

Replacement preserves semantic identity:

\[
\frac{
  \operatorname{id}(n')=i
  \quad
  \exists n\in N:\operatorname{id}(n)=i
  \quad
  \Gamma;((N-\{n\})\cup\{n'\},R)\vdash n'
}{
  (N,R)\xrightarrow{\operatorname{ReplaceNode}(i,n')}
  ((N-\{n\})\cup\{n'\},R)
}.
\]

Preserving the identifier does not automatically mean every semantic reference
remains valid. Evaluators and reference resolution determine the consequences
of the changed content.

\subsection{Remove nodes}

Removal can be defined without an implicit cascade:

\[
(N,R)\xrightarrow{\operatorname{RemoveNodes}(J)}
(N-\{n\in N\mid\operatorname{id}(n)\in J\},R-J).
\]

References from remaining nodes may subsequently become missing. A higher-level
cascade operation can be defined as a computed event containing all required
removals, but it should not be hidden inside the primitive rewrite.

\subsection{Example: atomic rectangle creation}

Creating a solver-native rectangle may add four lines and eight constraints.
If the fourth line has a duplicate identifier, applying only the first three
lines would leave an unintended partial rectangle. Atomic event semantics
require the complete creation event to fail, leaving the original snapshot
unchanged.

\section{Authored, Derived, and Diagnostic Facts}

\subsection{Motivation}

The specification places all facts in one dialect, which is useful, but it needs
rules preventing solver output from being confused with authored intent.

\subsection{Fact classes}

Partition the live nodes by semantic class:

\[
N=N_A\uplus N_D\uplus N_X,
\]

where:

\begin{itemize}
  \item $N_A$ contains authored or imported design intent,
  \item $N_D$ contains derived evaluator output, and
  \item $N_X$ contains diagnostics and observations.
\end{itemize}

This is a semantic partition. It may be represented by node type, provenance,
authority, or another mechanism.

Every derived or diagnostic node has provenance:

\[
\operatorname{prov}(n)=(a,v,Q,f),
\]

where:

\begin{itemize}
  \item $a\in A$ is the producing authority,
  \item $v$ identifies the evaluator semantics or version,
  \item $Q\subseteq I$ is the set of consulted source nodes, and
  \item $f$ is a fingerprint of the relevant source facts.
\end{itemize}

The core authority rule is:

\begin{quote}
An evaluator may create, replace, mark stale, or remove facts for which it has
authority. It may not silently rewrite authored intent.
\end{quote}

An evaluator may propose a repair to authored intent, but accepting that
proposal produces a separate authored event with an appropriate authority.

\subsection{Example: constraint repair}

A line is authored with endpoints $(0,0)$ and $(9.8,0.2)$, together with
horizontal and length-10 constraints. The solver derives a solved line from
$(0,0)$ to $(10,0)$.

The solved coordinates are a derived fact. They do not replace the authored
coordinates. If the user chooses to apply the solved position, that choice is a
new authored event rather than an unobservable solver mutation.

\section{Provenance, Currency, and Staleness}

\subsection{Motivation}

The specification refers to source hashes but does not define the semantic
condition that a hash is intended to represent.

\subsection{Definitions}

Let:

\[
\operatorname{fingerprint}_S(Q)
\]

be an abstract fingerprint of the source facts in $Q$ that are relevant to an
evaluator. The concrete fingerprint algorithm is outside the abstract
specification.

A derived node $d$ with provenance $(a,v,Q,f)$ is \textbf{current} in snapshot
$S=(N,R)$ when:

\[
\operatorname{current}_S(d)
\iff
Q\subseteq\operatorname{ids}(N)
\land
f=\operatorname{fingerprint}_S(Q).
\]

It is \textbf{stale} when all required sources still exist but their relevant
facts have changed:

\[
\operatorname{stale}_S(d)
\iff
Q\subseteq\operatorname{ids}(N)
\land
f\neq\operatorname{fingerprint}_S(Q).
\]

It is \textbf{orphaned} when at least one required source no longer exists:

\[
\operatorname{orphaned}_S(d)
\iff
Q\nsubseteq\operatorname{ids}(N).
\]

Separating stale from orphaned is useful: stale output may be recomputed from
available inputs, while orphaned output first requires repairing or replacing a
missing reference.

\subsection{Relevance requirement}

The fingerprint should cover only source facts relevant to the evaluator
result. Renaming a sketch should not necessarily stale its generated geometry,
while changing a constrained point should.

The abstract requirement is:

\[
S\equiv_{E,Q}S'
\implies
\operatorname{fingerprint}_S(Q)
=
\operatorname{fingerprint}_{S'}(Q),
\]

where $S\equiv_{E,Q}S'$ means evaluator $E$ observes the same relevant facts
in both snapshots.

\subsection{Example}

An extrusion result records the profile, distance parameter, direction, and
operation as its source facts. Changing the display color of the sketch does not
stale the body. Changing the extrusion distance from 10 mm to 12 mm does.

\section{Evaluators and Evaluation Closure}

\subsection{Motivation}

The execution cycle explains when evaluators run, but not what an evaluator is
or what it means for evaluation to be complete.

\subsection{Evaluator definition}

An evaluator is an abstract partial function:

\[
E_i:\mathcal{S}\rightharpoonup\Delta,
\]

where $\Delta$ is an event or finite sequence of events containing derived
facts, diagnostics, staleness marks, or removals under that evaluator's
authority.

An evaluator is \textbf{applicable} when its required inputs exist and satisfy
its input typing rules. An applicable evaluator may still produce a diagnostic
instead of its normal output.

Let $\mathcal{E}=\{E_1,\ldots,E_m\}$ be the evaluator set for a chosen target.
One evaluation step is:

\[
S_{k+1}=\operatorname{step}_{\mathcal{E}}(S_k).
\]

A snapshot has reached \textbf{evaluation closure} for target $t$ when no
applicable evaluator can add a newer required fact or replace a stale required
fact:

\[
\operatorname{closed}_{\mathcal{E},t}(S)
\iff
\operatorname{step}_{\mathcal{E},t}(S)=S.
\]

This equality is semantic. Implementations may replace equivalent caches or
bookkeeping without changing design state.

\subsection{Desired evaluator properties}

The specification should identify which of these properties are required and
which are desirable:

\begin{itemize}
  \item \textbf{Soundness:} emitted facts satisfy the semantic contract of their
    node types.
  \item \textbf{Authority confinement:} an evaluator changes only facts under
    its authority.
  \item \textbf{Determinism:} equivalent inputs produce semantically equivalent
    outputs.
  \item \textbf{Termination:} repeated evaluation reaches closure for supported
    finite inputs.
  \item \textbf{Confluence:} different valid evaluator schedules reach
    equivalent closed snapshots.
  \item \textbf{Locality:} changes outside declared source facts do not change
    evaluator output.
\end{itemize}

If strict determinism or confluence is not required for numerical evaluators,
the specification should define an equivalence relation or tolerance under
which results are considered the same.

\subsection{Example: failed extrusion}

A profile evaluator detects a closed region. The extrusion evaluator attempts
to create a body, but the geometry is self-intersecting. It emits a diagnostic
sourced from the extrusion and profile instead of a generated body.

Evaluation has still reached closure if the diagnostic is current and no
evaluator can produce a body from those inputs. Failure to create geometry is
not the same as failure of the graph evaluation mechanism.

\section{Constraint Semantics}

\subsection{Motivation}

The constraint vocabulary describes intended meanings in prose, but it does not
define the state over which constraints operate or distinguish checking from
solving.

\subsection{Geometric state}

For a sketch $s$, let:

\[
x_s\in X_s
\]

be its geometric state. Components of $x_s$ are the degrees of freedom exposed
by its entities, such as point coordinates, circle centers, and radii. An
authored sketch supplies an initial or preferred state:

\[
x_s^0\in X_s.
\]

Each constraint $c$ induces a residual function:

\[
F_c:X_s\rightarrow\mathbb{R}^{m_c}.
\]

The constraint is satisfied at $x$ under tolerance $\varepsilon_c$ when:

\[
\lVert F_c(x)\rVert_{W_c}\leq\varepsilon_c,
\]

where $W_c$ permits dimension-aware scaling of residual components.

\subsection{Example residuals}

For a line with endpoints $p_0=(x_0,y_0)$ and $p_1=(x_1,y_1)$, a horizontal
constraint has residual:

\[
F_{\mathrm{horizontal}}(x)=y_1-y_0.
\]

A vertical constraint has residual:

\[
F_{\mathrm{vertical}}(x)=x_1-x_0.
\]

A length constraint with target $d$ has residual:

\[
F_{\mathrm{length}}(x)=\lVert p_1-p_0\rVert-d.
\]

Coincidence of points $p$ and $q$ has residual:

\[
F_{\mathrm{coincident}}(x)=p-q.
\]

Perpendicular line directions $u$ and $v$ have residual:

\[
F_{\mathrm{perpendicular}}(x)=u\cdot v.
\]

These equations define satisfaction. They do not prescribe how a solver finds
$x$.

\subsection{Constraint policies}

The three policies can be defined as follows:

\begin{itemize}
  \item A \textbf{driving} constraint participates in determining a solved state
    $x_s^*$.
  \item A \textbf{checked} constraint is evaluated at $x_s^*$ but does not alter
    the admissible solution set.
  \item An \textbf{advisory} constraint reports its residual without making
    failure of the relation an evaluation failure.
\end{itemize}

A solved state for driving constraints $C_D$ satisfies:

\[
\forall c\in C_D:
\lVert F_c(x_s^*)\rVert_{W_c}\leq\varepsilon_c.
\]

When several solutions exist, the specification may define the preferred
result abstractly as one minimizing deviation from authored intent:

\[
x_s^*\in
\operatorname*{argmin}_{x\in X_s}
d(x,x_s^0)
\quad
\text{subject to all driving constraints.}
\]

The metric $d$ can remain evaluator-defined. Stating the stability objective
explains why small edits should normally lead to nearby solutions.

\subsection{Constraint statuses}

\begin{itemize}
  \item \Type{satisfied}: a solution exists and the residual is within tolerance.
  \item \Type{unsatisfied}: evaluated geometry violates the relation.
  \item \Type{undefined}: the relation cannot be evaluated, such as an angle
    involving a zero-length line.
  \item \Type{ambiguous}: references or solution selection do not determine the
    intended relation uniquely.
  \item \Type{overconstrained}: a subset of driving constraints has no common
    solution, or contains redundancy under a policy that treats redundancy as
    overconstraint.
  \item \Type{stale}: the recorded status was derived from source facts that
    have changed.
\end{itemize}

\Type{underconstrained} should also be considered as a sketch-solve status. It
means the solution set retains degrees of freedom:

\[
\dim(\mathcal{S})>0,
\]

where $\mathcal{S}\subseteq X_s$ is the set satisfying the driving constraints.

\subsection{Example: constrained rectangle}

A rectangle represented by four lines has eight endpoint points before
shared-corner constraints are considered. Coincident corner constraints connect
adjacent endpoints, horizontal and vertical constraints establish orientation,
and width and height constraints establish size. Even then, the rectangle
retains global translation unless a point or center is fixed. The solver should
report the remaining two translational degrees of freedom rather than treating
the sketch as invalid.

\section{Quantities, Dimensions, and Frames}

\subsection{Motivation}

The specification requires compatible units and inherited frames, but the
mathematical rules for compatibility and coordinate conversion are implicit.

\subsection{Dimensions and quantities}

Let $\mathcal{D}$ be a free abelian group of physical dimensions generated by
base dimensions such as length $L$, angle $\Theta$, mass $M$, and time $T$. A
quantity is:

\[
q=(v,d,u),
\]

where $v$ is its numeric value, $d\in\mathcal{D}$ is its dimension, and $u$ is
a unit representation for that dimension.

Addition and comparison require equal dimensions:

\[
d(q_1)=d(q_2).
\]

Multiplication combines dimensions:

\[
d(q_1q_2)=d(q_1)+d(q_2),
\]

and division subtracts them. Unit conversion changes the numeric representation
but not the physical quantity.

For example, 10 mm and 1 cm may be compared because both have dimension $L$ and
normalize to the same physical quantity. Adding 10 mm to 30 degrees is
ill-typed regardless of numeric representation.

\subsection{Frames}

Let every frame $f$ define an invertible transform into its parent frame:

\[
T_f\in\operatorname{Aff}(n),
\]

or, for rigid mechanical frames, $T_f\in SE(3)$. If frame ancestry connects
$f_a$ and $f_b$, coordinate conversion is defined by composition:

\[
p_b=T_b^{-1}T_a p_a.
\]

Geometric relations involving entities from different frames must either
convert them into a common frame or be undefined when no common frame
relationship exists.

For example, a PCB connector frame may be defined relative to a board. An
enclosure cutout may be defined relative to a part. Assembly instances supply
transforms that bring both into an assembly frame. A distance constraint between
the connector and cutout is meaningful only after composing those transforms
into a shared coordinate system.

\section{Validity Judgments}

\subsection{Motivation}

The word \emph{valid} can refer to graph integrity, reference resolution,
evaluator currency, constraint satisfaction, geometric validity, or
manufacturability. A single predicate obscures useful partial states.

\subsection{Independent predicates}

The graph is \textbf{well-formed} when it satisfies identity, root, typing, and
required-reference invariants:

\[
\operatorname{WellFormed}_{\Gamma}(S).
\]

A target is \textbf{resolvable} when its required references resolve with
accepted cardinalities:

\[
\operatorname{Resolvable}(S,t).
\]

A target is \textbf{current} when every required derived fact is current:

\[
\operatorname{Current}(S,t).
\]

A target is \textbf{evaluable} when applicable evaluators have sufficient
well-typed input:

\[
\operatorname{Evaluable}_{\mathcal{E}}(S,t).
\]

A target is \textbf{satisfied} under policy $P$ when its required constraints
satisfy that policy:

\[
\operatorname{Satisfied}_{P}(S,t).
\]

A generated body is \textbf{geometrically valid} according to the topology and
geometry contract:

\[
\operatorname{GeometryValid}(S,b).
\]

These predicates should not be collapsed. In particular:

\[
\operatorname{WellFormed}(S)
\not\!\!\implies
\operatorname{Satisfied}(S,t),
\]

and:

\[
\operatorname{WellFormed}(S)
\not\!\!\implies
\operatorname{GeometryValid}(S,b).
\]

\subsection{Example}

A graph contains a well-typed extrusion whose profile is self-intersecting. Its
references resolve, its distance has a length unit, and the evaluator has
emitted a current diagnostic explaining why no solid exists.

The graph is well-formed and evaluation is current. The extrusion has not
succeeded, and no geometrically valid body exists. This is an inspectable design
failure, not graph corruption.

\section{Complete Worked Example}

Consider a sketch containing four lines intended to form a 40 mm by 20 mm
rectangle, followed by a 10 mm extrusion.

\subsection{Authored nodes}

Let the authored node set contain:

\begin{itemize}
  \item frame $f_{xy}$,
  \item sketch $s$ on $f_{xy}$,
  \item lines $l_1,l_2,l_3,l_4$,
  \item four coincidence constraints $c_1,\ldots,c_4$,
  \item two horizontal constraints,
  \item two vertical constraints,
  \item width constraint $c_w=40\,\mathrm{mm}$,
  \item height constraint $c_h=20\,\mathrm{mm}$, and
  \item extrusion $x$ with distance $10\,\mathrm{mm}$.
\end{itemize}

These nodes are authored facts even if the initial line coordinates are
imperfect.

\subsection{Sketch evaluation}

The sketch evaluator reads the sketch entities and constraints:

\[
Q_s=\{s,l_1,l_2,l_3,l_4,c_1,\ldots,c_h\}.
\]

It finds a solved state $x_s^*$ satisfying the driving constraints and emits:

\begin{itemize}
  \item a \Type{SolvedSketch} node, and
  \item an underconstrained diagnostic reporting two global translational
    degrees of freedom.
\end{itemize}

The diagnostic does not make the graph malformed. It describes the solution
space.

\subsection{Profile and feature evaluation}

The profile evaluator consumes the solved sketch and identifies one bounded
region $p$. Its provenance records the sketch solution and relevant entity
identities. The profile is a derived fact, not authored topology.

The extrusion evaluator resolves its profile selector to $\{p\}$, verifies that
$10\,\mathrm{mm}$ has dimension $L$, and emits a generated body $b$ with
provenance from $x$ and $p$.

The dependency path is:

\begin{verbatim}
authored lines + constraints
            |
            v
       solved sketch
            |
            v
          profile
            |
            v
     generated body
\end{verbatim}

\subsection{Editing the width}

An authored event replaces $c_w$ with a constraint whose value is
$50\,\mathrm{mm}$ while preserving the constraint identifier. Immediately after
the event:

\begin{itemize}
  \item the authored graph remains well-formed,
  \item the old solved sketch is stale,
  \item the old profile is stale because its source solution is stale, and
  \item the generated body is stale through the transitive provenance chain.
\end{itemize}

Evaluation produces a new solved sketch, profile, and body. Stable authored
identifiers remain unchanged. Derived topology identifiers may be preserved
when semantic correspondence exists; otherwise dependent selectors become
missing or ambiguous rather than silently selecting unrelated topology.

\subsection{Removing a line}

If $l_3$ is removed without a cascade:

\begin{itemize}
  \item constraints referring to $l_3$ remain authored but have missing
    references,
  \item the sketch is no longer resolvable for profile evaluation,
  \item old derived results become orphaned or stale, and
  \item diagnostics identify the missing line references.
\end{itemize}

This state remains well-defined as a partially broken design. A later event may
restore a line with the same semantic identity or repair the affected
constraints.

\section{Suggested Organization of the Main Specification}

These additions can be incorporated without turning the main specification
into an implementation document. A possible structure is:

\begin{enumerate}
  \item \textbf{Mathematical preliminaries:} identifiers, types, references,
    snapshots, and histories.
  \item \textbf{Typed node graph:} type environment, well-formedness, edge
    roles, and roots.
  \item \textbf{Reference semantics:} set-valued resolution, cardinality,
    selectors, and stability.
  \item \textbf{Transition semantics:} atomic events and rewrite rules.
  \item \textbf{Fact classes and provenance:} authored, derived, diagnostic,
    current, stale, and orphaned.
  \item \textbf{Evaluation semantics:} evaluators, authority, closure, and
    diagnostics.
  \item \textbf{Domain relations:} constraint satisfaction, quantities, units,
    and frames.
  \item \textbf{Validity judgments:} structural validity separated from
    evaluation and geometry success.
  \item \textbf{Node vocabulary:} the existing domain-specific node tables.
  \item \textbf{Worked examples:} small normative examples exercising the
    abstract rules.
\end{enumerate}

The existing vocabulary can then be understood as a collection of type
declarations inside this semantic framework. The framework says what every
conforming PuppyCAD graph means; later documents may independently define
concrete syntax, storage, APIs, kernels, solvers, or user interaction.

\end{document}
